%──────────────────────────────────────────────────────────────────
\subsection{Resolution of the small-$\zeta$ regime: matched-NLSM
            and the $\sqrt{\zeta}$ scaling}
\label{subsec:matching}
%──────────────────────────────────────────────────────────────────

The Le~Gal--Schir\'o derivation gradient-expands the bosonic
action~\eqref{eq:bosonic-action} around the Lindbladian saddle of
Section~\ref{subsec:HS-saddle}, treating the dressed length
$\xi_{\rm dressed}\sim v_0/(\zeta\gamma)$ of
Remark~\ref{rem:dressed-length} as the short-distance UV cutoff.  This
prescription is well-adapted to the Born-rule endpoint $\zeta=1$ but
fails to describe the full $(\lambda,\zeta)$ phase diagram of
Case~B at small $\zeta$, where individual trajectories spend long
intervals under purely no-click evolution.  The resolution, developed
in this subsection, is a two-stage RG argument that runs the bare
cross-replica coupling on the free-fermion fixed point up to a
$\lambda$-controlled \emph{crossover} length scale $\xi_\lambda^{\rm
cross}$, and then matches to the DIII NLSM.  The output is a
parameter-free prediction $\lambda_c\propto\sqrt{\zeta}$ for the
location of the critical line at small $\zeta$, in quantitative
agreement with the cloning numerics of
Section~\ref{sec:binder-analysis}.

\subsubsection{The two-stage matched-NLSM argument}
\label{subsubsec:two-stage}

The matched argument is most cleanly phrased as a linearised RG flow
around the multicritical point
$(\lambda,\zeta)=(0,0)$, followed by the imposition of an order-one
threshold condition.

\paragraph{Multicritical fixed point and its relevant deformations.}
The point $(\lambda,\zeta)=(0,0)$ is a multicritical free-Dirac fixed
point of the replica Keldysh theory: at $\lambda=0$ the no-click
generator $H_{\rm nH}$ reduces to a real hopping Hamiltonian, and at
$\zeta=0$ the cross-replica vertex of Section~\ref{subsec:keldysh-action}
vanishes.  Both $\lambda$ and $\zeta$ are relevant
perturbations of this fixed point.  Linearising the RG equations
around the origin,
\begin{equation}\label{eq:RG-linear}
  \frac{d\lambda}{d\ell}\;=\; y_\lambda \,\lambda,
  \qquad
  \frac{d\zeta}{d\ell}\;=\; y_\zeta \,\zeta,
\end{equation}
the two eigenvalues are identified as follows.

\textbf{(i)~Cross-Choi eigenvalue $y_\zeta$.}~The cross-replica vertex
$\mathcal V_\zeta = \prod_r L_{x,r,+}L^\dagger_{x,r,-}$ is a product of
density bilinears (one per Keldysh contour) at coincident spacetime
points.  By engineering dimension counting at the free-Dirac UV with
$[\psi_\pm]=1/2$ and $[L_j]=1$ in $1+1$~d, its UV scaling dimension is
\begin{equation}\label{eq:Delta-zeta-UV}
  \Delta_\zeta^{\rm UV} \;=\; 1
  \qquad\Longrightarrow\qquad
  y_\zeta \;=\; 2-\Delta_\zeta^{\rm UV} \;=\; 1.
\end{equation}
This is confirmed by a direct lattice computation of the cross-Choi
two-point function $\langle\mathcal V_\zeta(x)\,\mathcal V_\zeta(0)\rangle\sim
x^{-2\Delta_\zeta^{\rm UV}}\sim x^{-2}$ at $\lambda=0$, presented in
Section~\ref{sec:cross-choi}.

\textbf{(ii)~$\lambda$-eigenvalue $y_\lambda$.}~The infrared fixed
point of the multicritical class-DIII NLSM is well established in the
disordered-fermion literature~\cite{KonigBrouwer2014} and was
transcribed to the MIPT context by~\cite{LeGalSchiro2025}; both give
a correlation-length exponent $\nu=2$.  Identifying $y_\lambda$ with
the inverse correlation-length exponent along the $\lambda$ direction
at the multicritical point,
\begin{equation}\label{eq:y-lambda}
  y_\lambda\;=\;\nu^{-1}\;=\;\tfrac{1}{2}.
\end{equation}
Equivalently, the $\lambda$ direction generates a crossover length
$\xi_\lambda^{\rm cross}\sim\lambda^{-1/y_\lambda}=\lambda^{-2}$.

\paragraph{Crossover length and matching condition.}
We emphasise the distinction between two correlation lengths that
share the same $\lambda$-scaling but are conceptually different.  The
\emph{crossover} length $\xi_\lambda^{\rm cross}$ is the RG scale at
which the linearised perturbation $\lambda(\ell)$ in
\eqref{eq:RG-linear} reaches order unity, signalling that the system
has flowed out of the vicinity of the $(0,0)$ multicritical fixed
point: $\lambda\,e^{y_\lambda\ell_\lambda}\sim 1$ gives
$e^{\ell_\lambda}\sim\lambda^{-2}$ and hence
\begin{equation}\label{eq:xi-lambda-cross}
  \xi_\lambda^{\rm cross}\;\sim\;\lambda^{-2}\,a.
\end{equation}
This is the length used in the matching condition below.  The
\emph{critical} length $\xi_{\rm crit}\sim|\lambda-\lambda_c(\zeta)|^{-\nu}$
is the diverging correlation length at the phase boundary itself and
plays no role in the present argument; using
$\xi_\lambda^{\rm cross}$ in place of $\xi_{\rm crit}$ avoids any
circularity in which the unknown $\lambda_c(\zeta)$ would be used to
fix itself.

Evaluating the linearised flow of $\zeta$ at the crossover scale,
\begin{equation}\label{eq:zeta-eff}
  \zeta_{\rm eff}(\xi_\lambda^{\rm cross})
  \;=\;
  \zeta\,e^{y_\zeta\ell_\lambda}
  \;=\;
  \zeta\,\lambda^{-y_\zeta/y_\lambda}
  \;=\;
  \zeta\,\lambda^{-2},
\end{equation}
where the last equality uses $y_\zeta/y_\lambda=2$.  The MIPT occurs
when this renormalised coupling reaches an order-unity threshold
$K^\ast$, the universal critical coupling of the class-DIII NLSM:
$\zeta_{\rm eff}(\xi_\lambda^{\rm cross})=K^\ast$, which solves to
\begin{equation}\label{eq:caseB-sqrtz}
  \boxed{\quad
  \lambda_c(\zeta)\;=\;\sqrt{c_\lambda/K^\ast}\,\sqrt{\zeta}
  \;\equiv\;C\,\sqrt{\zeta},
  \qquad
  C \;=\; O(1) \text{ non-universal}
  \quad}
\end{equation}
in the small-$\zeta$ regime $\zeta\ll 1$.  Here $c_\lambda$ is the
$O(1)$ coefficient relating the bare scale $a$ to the crossover scale
$\xi_\lambda^{\rm cross}=c_\lambda\,\lambda^{-2}\,a$, set by lattice
matching.

\paragraph{Two ingredients, two limits.}
The derivation rests on exactly two field-theory inputs:
$\Delta_\zeta^{\rm UV}=1$, established microscopically by the
cross-Choi lattice calculation; and $\nu=2$, the multicritical
class-DIII NLSM correlation-length exponent, taken from the literature
and consistent with the empirical $\nu(\zeta)$ extracted by global FSS
in the clean range $\zeta\in[0.1,0.7]$ (Section~\ref{sec:binder-analysis}).
The prefactor $C$ in \eqref{eq:caseB-sqrtz} is non-universal: it
depends on $c_\lambda$ and on the regularisation scheme defining
$K^\ast$, both of which are scheme-dependent $O(1)$ numbers.  Demanding
that the formula \eqref{eq:caseB-sqrtz} interpolate to the Born-rule
endpoint $\lambda_c(1)=1/2$ of~\cite{Carollo2018PRA} fixes the natural
one-parameter form
\begin{equation}\label{eq:lambda-c-full}
  \lambda_c(\zeta)\;=\;\frac{\sqrt{\zeta}}{1+\sqrt{\zeta}},
\end{equation}
equivalently $\alpha_c/w=\sqrt{\zeta}$, which is the closed-form
zero-parameter interpolation used throughout
Section~\ref{sec:binder-analysis}.

\subsubsection{Validity of the matched argument and the assumptions
               that go in}
\label{subsubsec:validity}

The derivation above is a \emph{matched crossover} argument, not a
complete one-loop calculation of the IR scaling dimensions of all
operators in the class-DIII NLSM.  We list explicitly the assumptions
on which the conclusion $\lambda_c\sim\sqrt{\zeta}$ rests and note,
for each, the level of justification provided.

\begin{itemize}
  \item \textbf{$\Delta_\zeta^{\rm UV}=1$ at the free-Dirac fixed point.}
    This is established by engineering dimension counting at the
    Keldysh-replica level and verified independently by lattice
    computation of the cross-Choi two-point function at $\lambda=0$,
    which decays as $x^{-2}$ (Section~\ref{sec:cross-choi}).  No
    further assumptions enter.

  \item \textbf{$\nu=2$ for the class-DIII multicritical NLSM.}
    This is a known result of the disordered-fermion NLSM
    literature~\cite{KonigBrouwer2014}, used in the MIPT context
    by~\cite{LeGalSchiro2025}.  Our global FSS gives values of
    $\nu(\zeta)$ scattered around $2$ in the clean range
    $\zeta\in[0.1,0.7]$, with finite-size systematics at the extremes
    preventing a precision determination from numerics alone.  We
    therefore use $\nu=2$ as a literature input, with the data
    serving as a consistency check.

  \item \textbf{The critical condition is reached at the crossover scale.}
    The matched argument runs $\zeta$ with the UV eigenvalue
    $y_\zeta^{\rm UV}=1$ from the lattice scale $a$ to the crossover
    scale $\xi_\lambda^{\rm cross}$, then sets the result equal to
    $K^\ast$.  This is justified \emph{provided} the IR running of
    $\zeta$ inside the NLSM regime is not parametrically large and
    does not change the leading exponent.  An anomalous IR dimension
    $\Delta_\zeta^{\rm IR}\neq 1$ would shift the prefactor and could
    generate subleading corrections; it would not change the leading
    $\sqrt{\zeta}$ behaviour unless the second RG window
    $\xi_\lambda^{\rm cross}\ll \ell\ll\xi_{\rm crit}$ is
    parametrically large compared with the first.  The full one-loop
    computation of $\Delta_\zeta^{\rm IR}$ at the class-DIII NLSM
    fixed point is left as future work.

  \item \textbf{No microscopic single-particle counterpart claimed.}
    Earlier formulations of this matching identified
    $\xi_\lambda^{\rm cross}$ with a BdG localisation length
    $\xi_{\rm ps}\sim\lambda^{-2}$ of the no-click Hamiltonian.  This
    identification is suggestive for models where the no-click
    dynamics has a manifest non-Hermitian SSH structure, but it does
    not directly map onto the bulk spectral content of
    $H_{\rm eff}=H-i\alpha\sum_j\tilde n_j$ for the
    Bogoliubov-mode-measurement model studied here, as direct
    diagonalisation shows.  We therefore treat
    $\xi_\lambda^{\rm cross}$ as the RG-defined crossover scale of
    Eq.~\eqref{eq:xi-lambda-cross}, with $\nu=2$ supplied by the
    field theory rather than by single-particle BdG considerations.
\end{itemize}

The conclusion is that the $\sqrt{\zeta}$ scaling of the Case~B
critical line is a robust universal prediction of the matched
field-theory argument, while the prefactor $C$ is non-universal and
not separately fixed by the present derivation.

\subsubsection{Comparison with numerics}
\label{subsubsec:numerics-match}

Global FSS on the merged Run-A, Run-B and Run-C cloning data
(Section~\ref{sec:binder-analysis}) gives, on the clean range
$\zeta\in[0.03,0.85]$,
\begin{align}
  \phi^{\rm fit} &\;=\; 0.56 \,\pm\, 0.05, \\
  C^{\rm fit}    &\;=\; 1.02 \,\pm\, 0.10,
\end{align}
with $\chi^2/\text{dof}=3.8$ for the free power-law fit
$\alpha_c/w=C\zeta^\phi$.  The prediction $\phi=1/2$ of
Eq.~\eqref{eq:caseB-sqrtz} is within $1.3\sigma$ of the best-fit
value; the alternative $\phi=1$, corresponding to a naive
$\xi_\lambda\sim\lambda^{-1}$ matching length, is excluded at
$9.1\sigma$.  The fitted prefactor is consistent with the order-unity
expectation of the matched argument.  The natural variable
$\alpha_c/w$ is approximately linear in $\sqrt{\zeta}$ with slope
close to one, reproducing the Carollo Born-rule endpoint
$\lambda_c(1)\to 1/2$ to within finite-size systematics
(Section~\ref{sec:binder-analysis}).

\paragraph{Crossover to the large-$\zeta$ regime.}
At $\zeta=O(1)$ the crossover length
$\xi_\lambda^{\rm cross}\sim\lambda^{-2}$ becomes comparable to the
lattice scale already at moderate $\lambda$, and the matched-NLSM
window shrinks.  In this regime the Lindbladian saddle of
Section~\ref{subsec:HS-saddle} is dominant and the full class-DIII
NLSM with bare coupling $g_B\propto\zeta\lambda/(1-\lambda)$ controls
the transition, recovering the Born-rule endpoint $\lambda_c\to 1/2$
at $\zeta=1$.  The closed-form $\lambda_c=\sqrt{\zeta}/(1+\sqrt{\zeta})$
of Eq.~\eqref{eq:lambda-c-full} interpolates between the two regimes
and is consistent with the numerics across the full range with no
additional fit parameters.

\begin{remark}[Physical picture]\label{rem:physical-interpretation}
The Case~B phase diagram is controlled by two structurally distinct
limits of the same DIII NLSM:
\begin{itemize}
  \item \textbf{Multicritical regime} ($\zeta\ll 1$, $\lambda\ll 1$):
    proximity to the free-Dirac fixed point at the origin.  The
    cross-replica vertex is a relevant operator of UV dimension
    $\Delta_\zeta^{\rm UV}=1$, and the linearised flow generates a
    $\lambda$-crossover length $\xi_\lambda^{\rm cross}\sim\lambda^{-2}$.
    Critical line $\lambda_c\propto\sqrt{\zeta}$, exponent
    $\nu=1/y_\lambda=2$.
  \item \textbf{NLSM regime} ($\zeta=O(1)$): the Lindbladian saddle
    is dominant and the DIII NLSM applies with bare coupling
    $g_B\propto\zeta\lambda/(1-\lambda)$.  Critical line approaches
    $\lambda_c=1/2$ at $\zeta=1$.
\end{itemize}
The two regimes are joined smoothly by the closed-form interpolation
\eqref{eq:lambda-c-full}; the same field-theoretic framework controls
both, with the UV cutoff in the gradient expansion taken to be the
lattice scale $a$ for $\zeta=O(1)$ and the crossover scale
$\xi_\lambda^{\rm cross}$ for $\zeta\ll 1$.
\end{remark}
